\section{Full rubric with anchors}
\label{app:rubric}

This appendix reproduces the frozen sixteen-parameter rubric in full. It is
the source-of-truth for every score in \texttt{scoring/scores.csv};
weights sum to $100$. Parameters 1--10 are the engineering-rigor
artefact-and-process dimensions; their relative weights are scaled so
that the ten of them sum to 70. Parameters 11--16 are the six
tool-surface dimensions (Section~\ref{subsec:rubric}) that prevent
the rubric from biasing toward tools that ship persistent multi-file
artefacts; each carries a flat weight of 5 (the six sum to 30). Each
parameter carries five integer anchors. When a run sits between two
anchors the scorer prefers the lower.

\subsection*{1.\ Artifact completeness --- weight 7}
Does the planner produce the \emph{shape} of artefacts a downstream
implementer needs (proposal, design, requirements, task list), or only
a chat blob?
\begin{description}[leftmargin=11mm, labelindent=0mm]
\item[1] A single chat-resident blob; no separable artefacts.
\item[2] A single markdown document, sectioned but not split.
\item[3] Two artefacts (e.g.\ a plan doc plus a task list) with light cross-linking.
\item[4] Three or more artefacts but missing one of \{proposal, design, requirements, tasks\}.
\item[5] Multi-file structured set covering proposal, design, requirements/specs, and tasks.
\end{description}

\subsection*{2.\ Requirement traceability --- weight 10.5}
Can a future reader trace each task back to a specific requirement?
\begin{description}[leftmargin=11mm, labelindent=0mm]
\item[1] Tasks have no link to requirements; ``implement X'' with no reference.
\item[2] Free-text references to requirements but no IDs.
\item[3] Some tasks cite requirement names but not all.
\item[4] Every task cites a requirement, but the format is inconsistent.
\item[5] Every task carries an explicit \texttt{Covers: <REQ-ID>} (or equivalent stable ID) and the IDs resolve to scenarios in the spec.
\end{description}

\subsection*{3.\ Task granularity --- weight 7}
Are tasks sized for individual review and check-off?
\begin{description}[leftmargin=11mm, labelindent=0mm]
\item[1] One mega-task (``implement the feature'') or no task list.
\item[2] A handful of multi-day chunks; not individually reviewable.
\item[3] $\sim 10$ medium tasks; some are still vague.
\item[4] $\sim 20$ tasks; most are concrete but a few are still mega-tasks.
\item[5] Tasks scoped to $\le 30$ minutes of work each; every task is a single PR-sized unit.
\end{description}

\subsection*{4.\ Test-plan coverage --- weight 10.5}
Does each task carry a test plan a reviewer can execute?
\begin{description}[leftmargin=11mm, labelindent=0mm]
\item[1] ``Add tests'' appears once globally with no specifics.
\item[2] A separate ``tests'' task at the end with no per-task linkage.
\item[3] Some tasks include a sentence-level test note.
\item[4] Most tasks include a concrete test plan, but coverage gaps remain.
\item[5] Every task ships with a per-task test plan (commands or scenario refs); test plans cite real fixtures or scenarios.
\end{description}

\subsection*{5.\ Code-reference accuracy --- weight 7}
Are file paths, function names, and APIs cited in the plan real, or hallucinated?
\begin{description}[leftmargin=11mm, labelindent=0mm]
\item[1] Referenced paths or APIs do not exist in the live repo.
\item[2] Some references are real, but at least one is hallucinated.
\item[3] All references are plausible; spot-check finds them mostly correct but with stale signatures.
\item[4] All file paths verified; minor naming drift in one or two function references.
\item[5] Every cited file/function/API is verifiable in the live repo at the planner's stated path.
\end{description}

\subsection*{6.\ MCP / external-source integration --- weight 7}
When the prompt references external systems (Jira ticket, GitHub PR), does the
planner actually read them?
\begin{description}[leftmargin=11mm, labelindent=0mm]
\item[1] Planner ignores the external link entirely.
\item[2] Planner names the external system but quotes nothing from it.
\item[3] Planner summarises the external source but with at least one factual error.
\item[4] Planner integrates the external source faithfully; minor omission.
\item[5] Planner reads, summarises, and cites the external system correctly; integrates ticket comments and linked code into the plan.
\end{description}

\subsection*{7.\ Reproducibility (artefact persistence) --- weight 7}
If the user closes the chat, can the plan be re-opened verbatim?
\begin{description}[leftmargin=11mm, labelindent=0mm]
\item[1] Plan lives only in volatile chat history; no on-disk artefact.
\item[2] User can manually copy the plan out, but no native export.
\item[3] Plan has a native on-disk surface, but no version-control integration.
\item[4] Plan is on disk and gits cleanly, but lacks a validation/lint step.
\item[5] Plan is a versioned directory in git with a CLI-driven validator (\texttt{openspec validate --strict} or equivalent).
\end{description}

\subsection*{8.\ Latency to ``ready to implement'' --- weight 3.5}
Wall-clock time from prompt-submitted to a complete plan being ready.
\begin{description}[leftmargin=11mm, labelindent=0mm]
\item[1] $> 60$ minutes.\quad
\item[2] $30$--$60$ minutes.\quad
\item[3] $15$--$30$ minutes.\quad
\item[4] $10$--$15$ minutes.\quad
\item[5] $< 10$ minutes.
\end{description}

\subsection*{9.\ Token / API cost --- weight 3.5}
Approximate API spend taken from \texttt{metadata.json.token\_estimate}.
\begin{description}[leftmargin=11mm, labelindent=0mm]
\item[1] $> \$5.00$.\quad
\item[2] $\$2.00$--$\$5.00$.\quad
\item[3] $\$1.00$--$\$2.00$.\quad
\item[4] $\$0.50$--$\$1.00$.\quad
\item[5] $< \$0.50$.
\end{description}

\subsection*{10.\ Hand-off friction to implementation --- weight 7}
Once the plan exists, how much friction is there to start coding from it?
\begin{description}[leftmargin=11mm, labelindent=0mm]
\item[1] Manual re-typing required; no path to feed the plan back to a coding agent.
\item[2] Plan can be copy-pasted into a fresh agent prompt, but no first-class hand-off tool.
\item[3] One-step manual hand-off (copy to file, point an agent at it).
\item[4] Native hand-off via a documented command, requiring one explicit user action.
\item[5] Single-command hand-off (\texttt{/opsx-apply}, Cursor $\to$ Agent Mode, Ralph loop) with no manual re-typing.
\end{description}

\subsection*{11.\ Tool licensing \& accessibility --- weight 5}
Is the tool itself open-source, and are the features that this study
exercises (MCP transports, agent mode, plan-card hand-off) reachable on
a free tier or behind a paid plan? Anchor evidence is drawn from the
public pricing pages~\citep{cline2026pricing,cursor2026pricing}.
\begin{description}[leftmargin=11mm, labelindent=0mm]
\item[1] Closed-source \emph{and} the primary feature exercised by the
study (chat MCP, agent mode) is gated behind a paid plan with no free
fallback.
\item[2] Closed-source freemium; the study's primary feature is gated
behind a paid plan (e.g.\ Cursor's MCP picker requires Pro at \$20/mo).
\item[3] Closed-source freemium where most features are reachable on
the free tier and only premium model rates require paid plans.
\item[4] Open-source frontend; some advanced features only available in
a paid hosted offering or with subscription credits.
\item[5] Fully open-source; all features used by this study reachable
on the free tier; BYOK supported (Cline; OpenSpec and Ralph as
free CLIs / loops with no licensing layer).
\end{description}

\subsection*{12.\ Learning curve / time-to-productivity --- weight 5}
How much friction does the operator face between ``never used the
tool'' and ``producing a usable plan''? This dimension is distinct
from O1 (setup friction) because it scores the cognitive cost of the
\emph{plan format and conventions} the operator must learn, not just
the install pipeline.
\begin{description}[leftmargin=11mm, labelindent=0mm]
\item[1] Multi-step CLI install plus multiple file conventions plus a
non-trivial mental model the operator must internalise (sprint
budgets, validators, lifecycle states) before the first plan.
\item[2] CLI install plus two coupled conventions to learn before any
plan content (OpenSpec: install + \texttt{openspec init} +
4-artefact mental model + \texttt{validate --strict} semantics).
\item[3] Single setup step plus one small mental model (Ralph:
\texttt{RALPH\_TASK.md} + loop iteration semantics).
\item[4] Already-in-IDE chat surface plus one micro-convention to
learn (e.g.\ a slash command or mode switch).
\item[5] Zero install, zero conventions; chat is the primary surface
and the tool drives the format end-to-end (Cursor in Plan Mode).
\end{description}

\subsection*{13.\ Human-in-the-loop validation experience --- weight 5}
How easily can the operator inspect, refine, and correct the plan
before it is handed to execution, without restarting the workflow?
This parameter is \emph{model-sensitive}: the underlying tool surface
sets the ceiling, but the model's habit of surfacing clarifying
questions versus committing to a single best guess shifts the score
within that ceiling.
\begin{description}[leftmargin=11mm, labelindent=0mm]
\item[1] No mid-flight refinement; any plan change requires restarting
the workflow from scratch with a new prompt.
\item[2] Plan can be edited but the tool must be re-invoked over the
edited file (a headless CLI re-run over an edited plan file:
edit-then-replay).
\item[3] Plan can be edited in place; the tool re-reads on next
invocation but offers no clarification affordance.
\item[4] Plan-card hand-off plus in-thread refinement; the tool
absorbs follow-up turns, but does not surface clarifying questions
to the operator (Cursor's chat surface; OpenSpec/Ralph under
Cline-as-extension --- both under either model).
\item[5] Plan-card hand-off plus in-thread refinement \emph{and} the
tool actively surfaces clarifying questions when the prompt is
ambiguous, raising the operator's awareness of the choices being
made (no cell in this study reaches anchor-5 on Param~13).
\end{description}

\subsection*{14.\ Autonomous capacity --- weight 5}
Once the plan exists, how much wall-clock work can the tool do without
the operator present, and how gracefully does it handle running out of
budget or hitting an unrecoverable failure?
\begin{description}[leftmargin=11mm, labelindent=0mm]
\item[1] No headless mode; an operator must be present to advance
every action and to dispatch every tool call (Cursor: even
``Agent Mode'' execution requires an attended chat panel).
\item[2] Headless mode exists, but with no budget intelligence and no
per-iteration commit (OpenSpec: end-to-end run drives the validator
to convergence but does not chunk progress to git).
\item[3] Headless mode plus per-iteration commit (each loop iteration
lands as a commit), but with no budget intelligence.
\item[4] Headless mode plus per-iter commit plus soft budget warnings
when wall-clock or token spend approaches a configured ceiling.
\item[5] Headless mode plus budget intelligence (cost ceilings,
sprint-style budgets), per-iter commit, and an explicit
\texttt{TASK\_GUTTER} fallback when the loop cannot make further
progress (Ralph: budgets, per-iter commit, and gutter fallback are
all built in).
\end{description}

\subsection*{15.\ Runtime integration \& ergonomics --- weight 5}
How well does the tool fit into the operator's existing IDE / shell
runtime, and how cleanly does it surface MCP transports and model
selection without requiring a separate process?
\begin{description}[leftmargin=11mm, labelindent=0mm]
\item[1] CLI scaffold install plus file-only conventions; no chat
surface and no integrated MCP runtime.
\item[2] CLI install with a separate runtime; chat surface available
only when paired with another tool (OpenSpec: paired with
Cline-as-extension or with \texttt{cursor-agent} for execution).
\item[3] Loop harness with optional in-IDE driver (Ralph: loop runs
on its own; Cline-as-extension and \texttt{cursor-agent} both
work as drivers).
\item[4] In-IDE chat panel with integrated MCP transports; runs
alongside the editor as a first-class extension.
\item[5] Native to the IDE the operator already uses; integrated MCP
picker; chat is the primary surface; model picker is one click;
slash-commands are first-class (Cursor).
\end{description}

\subsection*{16.\ Code-review / hand-off review surface --- weight 5}
Once a plan exists, how well can a non-author teammate (or a future
operator) review it asynchronously?
\begin{description}[leftmargin=11mm, labelindent=0mm]
\item[1] Chat thread only; no on-disk artefact a non-author can fetch
without operator help.
\item[2] Chat thread plus a non-deterministic local save location
(Cursor: \texttt{\textasciitilde{}/.cursor/plans/<hash>.plan.md}
exists locally but the path is operator-specific).
\item[3] Single file committed to git with no validator and no
required structure.
\item[4] Single file committed to git plus a non-strict structural
convention reviewers can read mechanically (Ralph:
\texttt{PLAN.md} sections + \texttt{TASK\_COMPLETE} sentinel).
\item[5] Multiple structured files committed to git plus a strict CLI
validator (\texttt{openspec validate --strict}) so reviewers can
audit completeness mechanically as well as visually.
\end{description}

\textbf{Frozen-as-of.} The rubric is frozen at the commit that
introduces \texttt{rubric.md}. Parameters 11--16 were added in the
revision that broadened the rubric to score tool-surface dimensions
(licensing, learning curve, human-in-the-loop validation, autonomous
capacity, runtime integration, and code-review surface). At
introduction, every prior row was re-scored under the new weights to
preserve the frozen-as-of contract: any change to parameters, weights,
or anchors after a single row has been written to \texttt{scores.csv}
requires re-scoring all prior rows.

\section{Operator-experience companion rubric}\label{app:opx-rubric}

The primary rubric in Appendix~\ref{app:rubric} is intentionally
weighted toward dimensions that matter for sustained, hand-off-ready
engineering work --- artefact persistence, requirement traceability,
and a strict validator. Lightweight surfaces such as a chat-resident
plan look bad on those dimensions even when they shine on the
operator's lived experience: setup time, cognitive load,
time-to-first-result, and how easy it is to course-correct without
restarting. To prevent the headline ranking from being read as a
property of the rubric instead of a property of the tools, this
appendix introduces a four-parameter \emph{operator-experience}
companion rubric with equal weights (25\,\% each, sum 100\,\%).

The companion rubric is scored \emph{per tool}, not per (tool $\times$
benchmark): O1--O4 are properties of the tool surface, not of the
task. Each parameter is scored on the same 1--5 anchor scale and
defended in \texttt{scoring/grill-me-log.md}.

\subsection*{O1.\ Setup friction --- weight 25}
How much work does the operator do before any plan content can
appear?
\begin{description}[leftmargin=11mm, labelindent=0mm]
\item[1] Multi-step CLI install plus scaffold conventions and file
layout the operator must learn before producing a usable plan.
\item[2] CLI install plus a small number of one-time configuration
files; conventions documented but still must be learned.
\item[3] CLI install only; conventions are minimal or self-documenting.
\item[4] No install; one configuration file (e.g.\ \texttt{.mcp.json})
is the entire setup.
\item[5] Zero install --- already in the IDE / shell the operator
uses for other work.
\end{description}

\subsection*{O2.\ Operator cognitive load --- weight 25}
How many sub-formats and conventions must the operator hold in head
to author and review a plan?
\begin{description}[leftmargin=11mm, labelindent=0mm]
\item[1] Multiple coupled sub-formats (proposal vs design vs spec
deltas vs tasks vs validator output); operator must remember which
content belongs in which file.
\item[2] Two-to-three coupled sub-formats with documented but
non-trivial separation.
\item[3] Single primary file plus one auxiliary surface (e.g.\
\texttt{PLAN.md} plus a task-completion signal).
\item[4] Single primary surface; auxiliaries are lightweight and
optional.
\item[5] Single chat surface; the tool drives the format and the
operator never has to recall a layout.
\end{description}

\subsection*{O3.\ Time-to-first-viable-plan --- weight 25}
Wall-clock time from prompt-submitted to the first piece of usable
plan content appearing on the operator's screen (not full plan
completion --- \emph{first} content).
\begin{description}[leftmargin=11mm, labelindent=0mm]
\item[1] Several minutes of scaffolding (file creation, validator
boot, MCP handshake) before any plan prose appears.
\item[2] Roughly a minute of scaffolding before plan prose starts
streaming.
\item[3] Some scaffolding overhead but the operator can see the plan
forming inside the first 30 seconds.
\item[4] Plan content streams within seconds; minimal scaffolding.
\item[5] Plan content streams immediately; the operator sees the
model thinking on screen.
\end{description}

\subsection*{O4.\ In-flight iterability --- weight 25}
How easy is it to refine the plan once it exists, without restarting
the workflow?
\begin{description}[leftmargin=11mm, labelindent=0mm]
\item[1] Plan revisions require restarting the workflow from scratch.
\item[2] Plan revisions require editing files by hand and re-running
the validator/loop.
\item[3] Plan revisions are supported via re-invocation against the
same files.
\item[4] Plan revisions are supported in-thread but require an
explicit ``revise'' command.
\item[5] Operator can refine the plan in the same chat without
restarting; the tool absorbs follow-up turns.
\end{description}

\subsection*{Per-tool scores}

The companion rubric is scored per tool (the four dimensions are
properties of the tool surface, not of the task). Anchors are written
defenses in \texttt{scoring/grill-me-log.md} (operator-experience
section). Table~\ref{tab:opx-scores} reports the per-parameter scores
and the equal-weighted per-tool totals.

\begin{table}[h]
\centering
\caption{Operator-experience scores (1--5 per parameter; weighted
total at equal 25\,\% weights, max 5.0). Column abbreviations: O1
setup friction, O2 cognitive load, O3 time-to-first-viable-plan, O4
in-flight iterability. OpenSpec and Ralph are scored under
Cline-as-extension, the driver used under \emph{both} models, so these
scores are model-stable. The OpenCode row is an anchor sketch
pending a formal grill-me rescoring pass; its native headless
operating mode constrains O4 to anchor-2 (revising a
plan requires re-invocation and overwriting the session), which caps
its OPX total well below Cursor's despite OpenCode leading the
engineering-rigor rubric.}
\label{tab:opx-scores}
\begin{tabular}{lcccccr}
\toprule
Tool & O1 & O2 & O3 & O4 & Total \\
\midrule
\texttt{cursor}     & 5 & 4 & 5 & 5 & \textbf{4.75} \\
\texttt{opencode} & 3 & 3 & 4 & 2 & \textbf{3.00} \\
\texttt{openspec}   & 2 & 2 & 3 & 3 & \textbf{2.50} \\
\texttt{ralph}      & 2 & 2 & 3 & 2 & \textbf{2.25} \\
\bottomrule
\end{tabular}
\end{table}

The ranking inverts the engineering-rigor headline: Cursor wins
the operator-experience rubric on the strength of zero-install setup
(O1: 5), single chat surface (O2: 4), immediate plan streaming (O3:
5), and in-thread revision (O4: 5). OpenCode places second: a
single-binary install plus a per-bench \texttt{opencode.json} is a
light ramp (O1: 3), \texttt{opencode run} streams content immediately
(O3: 4), and the operator holds only one plan file plus the
\texttt{session-export.json} schema in mind (O2: 3); its only real
weakness is that the headless-only operating mode admits revision only
by re-invocation and session overwrite (O4: 2). OpenSpec places third:
its 4-artefact scaffold imposes the highest file-format setup and
cognitive-load costs (O1: 2, O2: 2), though driven via
Cline-as-extension the agent thinking streams live in the chat panel
(O3: 3) and revisions are supported by re-invoking against the same
files (O4: 3). Ralph lands last on operator experience despite its
lightweight single \texttt{PLAN.md}: operating it requires authoring
\texttt{RALPH\_TASK.md} (success criteria plus a verification command)
and internalising the loop, budget, and
\texttt{GUTTER}/\texttt{COMPLETE} semantics --- the largest conceptual
ramp of the four (O1: 2, O2: 2) --- the loop boots before the first
plan prose appears (O3: 3), and refining the plan means editing the
task file and re-running the loop rather than taking a follow-up turn
in chat (O4: 2).
